\chapter{Seepage Coupling}

For the committed default benchmark, this chapter is intentionally short because
the seepage branch is \textbf{inactive}. The case file sets
\code{seepage = false}, so the mainline 3D PMG solve does not run any hydraulic
preprocessing or pore-pressure Newton loop.

\section{What The Mainline Mechanical Path Uses Instead}

On the default path:
\begin{itemize}
\item material properties come directly from the mechanical material table,
\item the body-force field comes directly from the mechanical
\code{gamma}-based volume load,
\item no pore-pressure field \code{pw} is solved,
\item no pressure gradient \code{grad\_p} is transferred into the mechanical
forcing,
\item no saturation update changes the material state before the mechanical
solve.
\end{itemize}

This is exactly why the current benchmark can go straight from mesh/material
preprocessing to constitutive setup and SSR continuation.

\section{Where The Branch Would Enter If Enabled}

If seepage were enabled, it would branch \emph{before} the mechanical load and
constitutive setup:
\[
\begin{aligned}
\text{mesh and material load} &\rightarrow \text{seepage solve} \\
&\rightarrow (\texttt{pw}, \texttt{grad\_p}, \texttt{mater\_sat}) \\
&\rightarrow \text{mechanical load and constitutive state}.
\end{aligned}
\]

That whole path is documented in Appendix~\ref{app:seepage}. It is real code,
but it is not part of the default parallel PMG run.
